\documentclass{amsart}
\usepackage{amssymb,mathtools}
\usepackage[hidelinks]{hyperref}
\hypersetup{pdftitle={The irrationality exponent of small-height lacunary series},pdfauthor={Ruiyi Zhang}}

\newtheorem{theorem}{Theorem}
\newtheorem{lemma}{Lemma}[section]
\numberwithin{equation}{section}
\newcommand{\Q}{\mathbb Q}
\newcommand{\Z}{\mathbb Z}
\newcommand{\R}{\mathbb R}
\newcommand{\den}{\operatorname{den}}

\begin{document}

\title[Small-height lacunary series]{The irrationality exponent of small-height lacunary series}
\author{Ruiyi Zhang}
\address{Independent Researcher, Columbia, Maryland, USA}
\email{ryzfsu@gmail.com}
\subjclass[2020]{Primary 11J82; Secondary 11J70, 11B83}
\keywords{irrationality exponent, lacunary series, rational height, continued fraction, folding}

\begin{abstract}
Let $N\ge2$ be an integer, let $u_1<u_2<\cdots$ be positive
integers with $\liminf u_{j+1}/u_j\ge2$, and let $\lambda_j$ be
nonzero rational numbers with $\log H(\lambda_j)=o(u_j)$.
Then $x=\sum_{j\ge1}\lambda_jN^{-u_j}$ is irrational and its
irrationality exponent is $\limsup u_{j+1}/u_j$, finite or infinite.
The proof compares logarithmic gaps between the continued-fraction
denominators of consecutive truncations. A perturbation lemma places
each gap of the next truncation in at least one of five classes (bounded,
inherited, central, short, reflected), each with an explicit error. Reflection
moves a large gap far enough to the right to give a uniform bound at
the boundary $\liminf u_{j+1}/u_j=2$.
The coefficient condition allows arbitrary signs and cannot be
replaced by $\log H(\lambda_j)\le C_*u_j$ for any fixed $C_*>0$.
\end{abstract}

\maketitle

\section{Introduction}

For an irrational real number $\xi$, its irrationality exponent is
\[
\mu(\xi)=\sup\left\{\nu\in\R:
\left|\xi-\frac pq\right|<q^{-\nu}
\text{ for infinitely many }(p,q)\in\Z\times\Z_{>0}\right\}.
\]
For a nonzero rational number written in lowest terms, put
\begin{equation}\label{eq:height}
H(a/b)=\max(|a|,b),\qquad a\in\Z\setminus\{0\},\quad
b\in\Z_{>0},\quad \gcd(a,b)=1.
\end{equation}
For $y\in\Q$, let $\den(y)$ denote its positive reduced denominator.
All logarithms are natural.

\begin{theorem}\label{thm:main}
Let $N\ge2$ be an integer, let $u_1<u_2<\cdots$ be positive
integers, and let $\lambda_j\in\Q\setminus\{0\}$ satisfy
\[
\log H(\lambda_j)=o(u_j),
\]
and suppose that
\[
\liminf_{j\to\infty}\frac{u_{j+1}}{u_j}\ge2.
\]
Then the absolutely convergent series
\[
x=\sum_{j\ge1}\lambda_jN^{-u_j}
\]
is irrational and
\begin{equation}\label{eq:main}
\mu(x)=\limsup_{j\to\infty}\frac{u_{j+1}}{u_j}.
\end{equation}
Both sides lie in $[2,+\infty]$.
\end{theorem}

The ratios need not converge. For example, integer exponent sequences
with repeating ratio patterns $(2,2,3)$ and $(2,5/2)$ give exponents
$3$ and $5/2$, respectively. Ratios may fall below $2$ by a quantity
tending to zero. The coefficients may have either sign; any sequence
from a finite subset of $\Q\setminus\{0\}$ is allowed. The binary specialization
$u_j=2^{j-1}$ of Theorem~\ref{thm:main} is the main theorem of
\cite{Zh}, proved there with Hankel determinants and Pad\'e approximants;
the present proof uses a different argument based on regular continued
fractions, and the sharpness examples of Section~9 follow Section~6 of
that paper.

For an integer base $b\ge2$ and a strictly increasing sequence of
positive integers $(n_j)$ with $n_{j+1}/n_j\ge2$ for all large $j$, the formula
$\mu(\sum_j b^{-n_j})=\limsup n_{j+1}/n_j$ for unit coefficients is a
classical consequence of folding, recorded in Bugeaud, Krieger and Shallit
\cite{BKS} with attribution there to Bugeaud \cite{Bu08}. Bugeaud \cite[Theorem~7.21]{Bu12} gives a result on prescribed
irrationality exponents for a specific family in the middle-third Cantor set.
The explicit Cantor-set family in Bugeaud \cite[Theorem~2]{Bu08} consists,
up to a fixed rational factor, of special cases of the unit-weight part of
Theorem~\ref{thm:main}.
Theorem~\ref{thm:main} also treats small rational weights and permits ratios
tending to $2$ from below.
Continued-fraction expansions of lacunary series appear in Shallit
\cite{Sh79,Sh82}. The reversal used here is
related to the folding identities of Mend\`es France \cite{MF}
and van der Poorten and Shallit \cite{vPS}.

Duverney, Kurosawa and Shiokawa \cite[Theorem~1]{DKS} give an
irrationality-exponent formula for series
$\sum_j\sigma_j/X_j$, where $\sigma_j\in\{-1,1\}$,
$X_j>1$ are rational, and $X_{j+1}\ge X_j^2$ eventually. Write $T_j$
for their auxiliary sequence of positive integers, defined by
\begin{equation}\label{eq:DKS}
T_1=\den(X_1),\qquad
T_{j+1}=T_j^2\den(X_{j+1}/X_j^2).
\end{equation}
The additional hypothesis is
$\log T_{j+1}=o(\log X_j)$.
Their conclusion is
\begin{equation}\label{eq:DKSconclusion}
\mu\!\left(\sum_j\frac{\sigma_j}{X_j}\right)
=\limsup_j\frac{\log X_{j+1}}{\log X_j}.
\end{equation}
For $\lambda_j\in\{-1,1\}$ and $u_{j+1}\ge2u_j$ eventually, discard
a finite prefix and take $X_j=N^{u_j}$. Then $X_{j+1}/X_j^2$ is an
integer and $T_j=1$, so the DKS formula gives \eqref{eq:main}, including
the case of exact ratio $2$.

For the following comparison with arbitrary rational weights, assume\linebreak
$\limsup u_{j+1}/u_j<\infty$, so that $u_{j+1}=O(u_j)$. In the subcase
$\liminf u_{j+1}/u_j>2$, discard finitely many terms and take
$X_j=N^{u_j}/|\lambda_j|$ and $\sigma_j=\operatorname{sgn}(\lambda_j)$.
The height hypothesis makes $\log X_j$ asymptotic to $u_j\log N$, so
$X_j>1$ eventually and the logarithmic ratios have the same limsup as the
exponent ratios. The strict lower-ratio margin then gives
$X_{j+1}\ge X_j^2$ eventually. The local logarithmic denominator costs are
$o(u_j)$: if
$\lambda_j=\alpha_j/\beta_j$ in lowest terms, then
\[
\den(X_{j+1}/X_j^2)\le |\alpha_{j+1}|\beta_j^2.
\]
Put $h_j=\log H(\lambda_j)/u_j$, so $h_j\to0$.
Since $u_{j+1}\ge(2+\varepsilon_*)u_j$ eventually
for some $\varepsilon_*>0$, the fixed initial denominator cost is negligible
after normalization, and the normalized later costs are bounded by a
convergent geometric sum. Their accumulation with weights $2^{j-i}$, $i\le j$, is
therefore still $o(u_j)$, giving $\log T_{j+1}=o(\log X_j)$ as required by DKS.
Adding back the omitted rational prefix does not change the irrationality
exponent. The case of an infinite limsup follows from the truncations alone
(Section~\ref{sec:arithmetic}).
Amou and Bugeaud \cite[Lemma~3]{AB} give a related criterion.
In notation local to this paragraph, let $\widehat A_n$ be strictly increasing
positive integers, let $\widehat B_n$ be nonzero integers, and let
$\widehat\tau=\limsup\log \widehat A_{n+1}/\log \widehat A_n$ be finite and at least two.
Their local condition is
\[
\frac{\widehat A_{n+1}/|\widehat B_{n+1}|}{(\widehat A_n/|\widehat B_n|)^2}\ge2
\quad\text{eventually},
\]
where $\widehat S_n=(\widehat A_{n+1}/|\widehat B_{n+1}|)/(\widehat A_n/|\widehat B_n|)^2$
denotes the left-hand side. Let
$\widehat d_n=\den(\widehat S_n)$, $\widehat d_0=|\widehat B_1|$, and
\[
\widehat D_n=\prod_{i=0}^n \widehat d_i^{2^{n-i}}.
\]
Their remaining conditions are $|\widehat B_n|=\widehat A_n^{o(1)}$ and
$\widehat D_n=\widehat A_n^{o(1)}$. Under these conditions, the irrationality exponent of
the series with terms $\widehat B_n/\widehat A_n$ equals $\widehat\tau$.
In the proof of their Theorem~2, Amou and Bugeaud \cite[Section~5]{AB}
state that there exists a sequence of ratios greater than $2$ with lower
limit $2$ for which Lemma~3 applies to their construction; the omitted
details concern the cumulative denominator condition
$\widehat D_n=\widehat A_n^{o(1)}$.

At the boundary $\liminf u_{j+1}/u_j=2$, for the choice
$X_j=N^{u_j}/|\lambda_j|$ above, the height hypothesis alone
does not imply $\log T_{j+1}=o(\log X_j)$: weighting local costs by
$2^{j-i}$ can leave a nonnegligible contribution. Also, with this choice,
a step $u_{j+1}<2u_j$ gives $X_{j+1}<X_j^2$ for unit coefficients.
Theorem~\ref{thm:main} includes arbitrary rational weights of the stated
height at this boundary and permits such steps below $2$.

The proof uses only regular continued fractions and rational arithmetic. The
continued fraction of one
truncation need not be a prefix of that of the next, so a perturbation
lemma controls each gap separately. The reflection estimate gives a lower
bound for the rightward displacement of a large gap in terms of its old
width. At the boundary ratio $2$ there is no margin, and this shift takes
its place. For the infinite continued fraction, each gap gets its own
truncation, chosen so that the accumulated finite-stage error stays
logarithmic.

\section{Continued-fraction facts}

For a nonintegral rational number we use its canonical finite regular
continued fraction, whose last partial quotient is at least $2$. For an
integer, the canonical expansion consists of that integer alone.
For irrational numbers we use infinite regular continued fractions.
Convergents are indexed from zero and written with positive denominators;
the auxiliary pair with denominator zero is not counted as a convergent.
A convergent preceding the terminal rational is called nonterminal.

If $p/q$ is a nonterminal convergent of a real number $\theta$ and
the next convergent has denominator $q^+$, then
\begin{equation}\label{eq:CFerror}
\frac1{q(q^++q)}\le\left|\theta-\frac pq\right|
\le\frac1{qq^+}.
\end{equation}
They also hold for rational $\theta$ with its canonical finite
expansion. We use Legendre's criterion in both cases:
if $p/q$ is reduced and
\begin{equation}\label{eq:Legendre}
\left|\theta-\frac pq\right|<\frac1{2q^2},
\end{equation}
then $p/q$ is a convergent of $\theta$. If $\theta$ is rational and
the error is nonzero, it is nonterminal. For a nonintegral rational $\theta$, let $q_{m-1}$ and $q_m$ be the
penultimate and final denominators in its canonical expansion.
The strict inequality rules out the penultimate convergent of the alternative
expansion ending in $1$: its denominator is $q_m-q_{m-1}$ and
its error is $1/(q_m(q_m-q_{m-1}))$, so the strict threshold cannot
hold when the canonical last partial quotient is at least $2$.
For integral $\theta$, a distinct rational with positive denominator has
error at least the reciprocal of that denominator, so the strict criterion
can hold only at $\theta$ itself. For irrational $\theta$ every convergent
is nonterminal.

A \emph{logarithmic gap} is a pair
$a=\log q$, $b=\log q^+$ from consecutive convergent denominators;
its width is $g=b-a$. A possible initial equality $q=q^+=1$ is
allowed.

\section{Reversing a finite denominator profile}

Let
\[
\alpha=\frac PV=[a_0;a_1,\ldots,a_m],\qquad V\ge2,
\]
be canonical. Write its convergent denominators as
$q_{-1}=0,q_0=1,\ldots,q_m=V$, and put
\[
V'=q_{m-1},\qquad
w_0=\frac{V'}V=[0;a_m,a_{m-1},\ldots,a_1],\qquad M=\log V.
\]

\begin{lemma}\label{lem:reversal}
For every logarithmic gap $\log s\to\log S$ in the canonical
finite continued fraction of $w_0$, there is a logarithmic gap
$c\to d$ of $\alpha$ such that
\begin{equation}\label{eq:reversal}
|\log s-(M-d)|\le\log2,\qquad
|\log S-(M-c)|\le\log2.
\end{equation}
In particular, $0\le c\le d\le M$.
\end{lemma}

\begin{proof}
Let $r_i$, $0\le i\le m$, be the denominators in the displayed
reversed expansion before a possible canonical merge, which is treated at
the end. Thus $r_0=1$
and $r_m=V$. Continuant symmetry and the splitting identity give
\[
V=q_jr_{m-j}+q_{j-1}r_{m-j-1}\qquad(0\le j\le m),
\]
where $r_{-1}=0$. Both denominator sequences are nondecreasing, so
\begin{equation}\label{eq:continuantbounds}
\frac{V}{2q_j}\le r_{m-j}\le\frac V{q_j}.
\end{equation}
Hence the gap $r_i\to r_{i+1}$ corresponds to
$q_{m-i-1}\to q_{m-i}$, with the errors in \eqref{eq:reversal}.

If $a_1\ge2$, the reversed expansion is canonical. If $a_1=1$,
its canonical form merges the final $1$ and removes only $r_{m-1}$.
The two final gaps merge into $r_{m-2}\to r_m=V$, which corresponds to
$q_1\to q_2$; all earlier gaps retain the preceding correspondence. Since
$q_1=1$, its upper endpoint
is exactly $M-\log q_1$, and its lower endpoint satisfies
\eqref{eq:continuantbounds} with $j=2$. The case $m=1,a_1=1$
is excluded by $V\ge2$.
\end{proof}

\section{A rational perturbation lemma}

Let $\alpha=P/V$ be reduced with $V\ge2$, and let
$\beta=\alpha+\delta$ with $\delta\in\Q\setminus\{0\}$.
Put
\begin{equation}\label{eq:perturbparameters}
\Delta=|\delta|,\qquad M=\log V,\qquad L=-\log\Delta,
\qquad W=\frac1{V^2\Delta}=e^{L-2M}.
\end{equation}
Let $P'/V'$ be the penultimate convergent in the canonical expansion
of $\alpha$, and write
\begin{equation}\label{eq:inversecoordinate}
\sigma=PV'-P'V\in\{1,-1\},\qquad
\zeta=-\frac{\sigma}{V^2\delta}.
\end{equation}
Put $B=\den(\zeta)$ and $E=\log(4B)$. So $B\ge1$, $|\zeta|=W$, and
$\zeta$ may be negative.

\begin{lemma}\label{lem:perturbation}
Let $p/q$ be a nonterminal convergent of $\beta$ with next denominator
$q^+$. Put
\[
\varepsilon=|\beta-p/q|>0,\qquad \pi_q=q^+/q,\qquad
a=\log q,\quad b=\log q^+,\quad g=b-a.
\]
At least one of the following alternatives holds.
\begin{enumerate}
\item \textbf{Bounded.} $\pi_q\le6$, so $g\le\log6$.
\item \textbf{Inherited.} $p/q\ne\alpha$ is a nonterminal convergent of $\alpha$ whose next
denominator $U$ satisfies $U\ge q^+/6$. In particular,
the new gap has the same left endpoint $a=\log q$, and
\begin{equation}\label{eq:inherited}
g\le\log U-\log q+\log6.
\end{equation}
\item \textbf{Central.} $p/q=\alpha$, so $q=V$ and
\begin{equation}\label{eq:central}
g\le L-2M.
\end{equation}
\item \textbf{Short.} One has
\begin{equation}\label{eq:short}
g\le\log12+2\log B\le2E,\qquad
a\ge L-M-\log(3/2).
\end{equation}
\item \textbf{Reflected.} There is an old logarithmic gap $c\to d$ of $\alpha$ such that
\begin{align}
|a-(L-d)|&\le E,\qquad |b-(L-c)|\le E,\label{eq:reflected}\\
\intertext{and}
a&\ge L-M-\log(3/2).\label{eq:reflectedlocation}
\end{align}
\end{enumerate}
No lower bound on $W$ is assumed.
\end{lemma}

\begin{proof}
Suppose $\pi_q>6$. If $\varepsilon\ge\Delta/2$, then
\begin{equation}\label{eq:oldlegendre}
\left|\alpha-\frac pq\right|
\le\varepsilon+\Delta\le3\varepsilon
\le\frac3{qq^+}<\frac1{2q^2}.
\end{equation}
If $p/q=\alpha$, then $q=V$ and $\varepsilon=\Delta$.
Equation~\eqref{eq:CFerror} gives $q^+\le1/(V\Delta)=VW$,
which is the central alternative. Otherwise, $p/q$ is reduced and differs
from $\alpha$, so Legendre's criterion gives a nonterminal convergent of
$\alpha$. Its next denominator $U$ satisfies
\[
\frac1{q(U+q)}\le\frac3{qq^+},\qquad
U\ge q^+/3-q\ge q^+/6,
\]
giving the inherited alternative.

Now suppose $\varepsilon<\Delta/2$. With $w_0=V'/V$, define the local coordinate $t=\zeta-w_0$. Direct calculation gives
\begin{equation}\label{eq:mobius}
\beta=\frac{Pt+P'}{Vt+V'},\qquad Vt+V'=V\zeta.
\end{equation}
Signs of $\zeta$ and $t$ are arbitrary.
For $f(w)=(Pw+P')/(Vw+V')$, the M\"obius difference formula is
\[
f(t)-f(w)=\frac{\sigma(t-w)}{(Vt+V')(Vw+V')}.
\]
Define integers $h,k$ by
\begin{equation}\label{eq:inversepair}
\binom pq=\begin{pmatrix}P&P'\\V&V'\end{pmatrix}\binom hk
\end{equation}
The matrix is unimodular, so $h$ and $k$ are coprime. Here $k\ne0$:
otherwise \eqref{eq:inversepair} would give $p/q=\alpha$, contrary to
the present error assumption. Moreover,
\[
k=\sigma(Pq-Vp),\qquad
K:=|k|=Vq\left|\alpha-\frac pq\right|.
\]
The second row of \eqref{eq:inversepair} gives $V(h/k)+V'=q/k$.
Substituting $w=h/k$ in the M\"obius difference formula gives an exact
expression for $|t-h/k|$, valid for either sign of $k$ and $\zeta$.
The assumption $\varepsilon<\Delta/2$ gives
\begin{equation}\label{eq:Kbounds}
\frac{q}{2VW}\le K\le\frac{3q}{2VW}.
\end{equation}
Since $K$ is a positive integer, $K\ge1$, so
\begin{equation}\label{eq:qafter}
q\ge\frac23VW.
\end{equation}
This is the bound on $a$ in the short and reflected alternatives.

Write $r=\zeta-h/k=y/s$ in reduced form with $s>0$.
Comparing the denominators of $\zeta$, $h/k$ and $r$ gives
\begin{equation}\label{eq:sbound}
K/B\le s\le BK.
\end{equation}
The bound $s\le BK$ comes from the definition of $r$, and
$K\le Bs$ from $h/k=\zeta-r$; signs do not matter.
The unimodular difference identity from
\eqref{eq:mobius}--\eqref{eq:inversepair} gives the local error $\omega_0$:
\begin{equation}\label{eq:inverseerror}
\omega_0:=|w_0-r|=\left|t-\frac hk\right|
=\frac{\varepsilon qVW}{K}>0.
\end{equation}
If $\pi_q\le12B^2$, the short alternative holds. Otherwise,
\begin{equation}\label{eq:inverselegendre}
\omega_0s^2\le B^2\varepsilon qVWK
\le\frac32B^2\varepsilon q^2
\le\frac{3B^2}{2\pi_q}<\frac18.
\end{equation}
Legendre's criterion makes $r$ a nonterminal convergent of $w_0$.
Let $S$ be its next denominator. Equation~\eqref{eq:CFerror} yields
\[
\frac1{s\omega_0}-s\le S\le\frac1{s\omega_0},\qquad
\frac1{2s\omega_0}\le S\le\frac1{s\omega_0},
\]
where the second lower bound uses $\omega_0s^2<1/8$.
Equations~\eqref{eq:CFerror} and
\eqref{eq:Kbounds}--\eqref{eq:inverseerror} now give
\begin{equation}\label{eq:reflecteddenominators}
\frac{q}{2BVW}\le s\le\frac{2Bq}{VW},\qquad
\frac{q^+}{2BVW}\le S\le\frac{2Bq^+}{VW}.
\end{equation}
For the second pair, use
\[
\frac1{s\omega_0}=\frac{K}{s\varepsilon qVW},\qquad
1/B\le K/s\le B,\qquad
q^+\le\frac1{\varepsilon q}\le q^++q.
\]
Since $\log(VW)=L-M$, \eqref{eq:reflecteddenominators} implies
\begin{equation}\label{eq:shiftedprofile}
\begin{aligned}
|\log s-[a-(L-M)]|&\le\log(2B),\\
|\log S-[b-(L-M)]|&\le\log(2B).
\end{aligned}
\end{equation}
Apply Lemma~\ref{lem:reversal} to the gap $s\to S$.
The combined endpoint error is at most
$\log(2B)+\log2=E$, giving \eqref{eq:reflected}.
\end{proof}

For a reflected gap with old width $g_0=d-c$,
\eqref{eq:reflected} gives
\begin{equation}\label{eq:widthandgain}
g\le g_0+2E,\qquad a-c\ge(L-2M)+g_0-E.
\end{equation}
For the second, use $d\le M$ and
$a-c\ge L-d-c-E=L-2d+g_0-E$.
When $L=2M+o(M)$, the second inequality bounds the rightward shift
$a-c$ from below by the old width $g_0$, up to the error terms.

\section{Arithmetic of the truncations}\label{sec:arithmetic}

Write $\lambda_j=\alpha_j/\beta_j$ in reduced form with $\beta_j>0$,
and put
\begin{equation}\label{eq:truncations}
t_n=u_n\log N,\qquad
x_n=\sum_{j\le n}\lambda_jN^{-u_j}=\frac{P_n}{V_n}
\end{equation}
in reduced form with $V_n>0$. The ratio assumption gives a fixed lower
ratio $\rho>1$ for all large indices, for example $\rho=3/2$.
The height hypothesis and geometric growth imply
\begin{equation}\label{eq:heightaccumulation}
\sum_{j\le n}\log H(\lambda_j)=o(u_n).
\end{equation}
Fix $\eta>0$. After an index depending on $\eta$, each summand is at
most $\eta u_j$, and
$u_j\le u_n\rho^{-(n-j)}$ there. Sum this geometric bound
and let $\eta$ tend to zero.

Set
\[
D_n=\prod_{j\le n}\beta_j,\qquad
S_n=\sum_{j\le n}\alpha_j(D_n/\beta_j)N^{u_n-u_j}.
\]
Then $x_n=S_n/(N^{u_n}D_n)$ and $\log D_n=o(u_n)$.
For every prime $p\mid N$, each summand of $S_n$ before the last
is divisible by $p^{v_p(N)(u_n-u_{n-1})}$. The last is
$\alpha_nD_{n-1}$,
whose valuation is $o(u_n)$ by \eqref{eq:heightaccumulation}.
Since $u_n-u_{n-1}\ge(1-1/\rho)u_n$ eventually, the last summand
has strictly smaller valuation. Hence eventually
\begin{equation}\label{eq:valuation}
v_p(S_n)=v_p(\alpha_nD_{n-1})=o(u_n).
\end{equation}
Since $N$ has only finitely many prime divisors, the valuation estimates
hold simultaneously after increasing the initial index. So $S_n\ne0$.
If $G_n=\gcd(S_n,N^{u_n}D_n)$,
the part supported outside the primes dividing $N$ is at most $D_n$,
and \eqref{eq:valuation} bounds the remaining part by $\exp(o(u_n))$.
Consequently,
\begin{equation}\label{eq:reduceddenominator}
\log G_n=o(u_n),\qquad
V_n=\frac{N^{u_n}D_n}{G_n},\qquad
M_n:=\log V_n=t_n+o(t_n).
\end{equation}

The first omitted term dominates the rest of the tail, giving
\begin{equation}\label{eq:tail}
x-x_n=\lambda_{n+1}N^{-u_{n+1}}(1+o(1)),\qquad
-\log|x-x_n|=t_{n+1}+o(t_{n+1}).
\end{equation}
For this and for absolute convergence, choose
$0<\tau<\log N\,(\rho-1)/(\rho+1)$ and use
$e^{-\tau u_j}\le|\lambda_j|\le e^{\tau u_j}$ at all late
indices. For every $k\ge n+1$ the ratio of successive absolute terms
is at most
\[
\exp\!\left(-[\rho(\log N-\tau)-(\log N+\tau)]u_k\right).
\]
So the rest of the tail is exponentially small relative to the first
omitted term, and the tail is eventually nonzero.

If $x=a'/b'$ were rational, its nonzero difference from $x_n$
would be at least $1/(b'V_n)$, contrary to
\eqref{eq:reduceddenominator}--\eqref{eq:tail} and the eventual
lower ratio $\rho>1$. Thus $x$ is irrational, and its truncations give
\begin{equation}\label{eq:lowerbound}
\mu(x)\ge\limsup_{n\to\infty}\frac{u_{n+1}}{u_n}.
\end{equation}
By \eqref{eq:reduceddenominator}, $V_n\to\infty$, so the truncations
give infinitely many distinct rational approximants. Their approximation exponents are
$t_{n+1}(1+o(1))/(t_n(1+o(1)))$. An infinite limsup gives
\eqref{eq:main} at once.

For the rest of the proof assume
\begin{equation}\label{eq:finiteR}
F:=\limsup_{n\to\infty}\frac{u_{n+1}}{u_n}<\infty.
\end{equation}
Thus $F\ge2$ and $u_{n+1}=O(u_n)$. Apply
Lemma~\ref{lem:perturbation} to
\[
\delta_n=x_{n+1}-x_n=\lambda_{n+1}N^{-u_{n+1}},\qquad
L_n=-\log|\delta_n|.
\]
Here the application has $\alpha=x_n$, $\beta=x_{n+1}$, $\delta=\delta_n$, so that $M=M_n$, $L=L_n$, $\zeta=\zeta_n$ and $B=B_n$.
The absolute value of the signed coordinate $\zeta_n$ is
\[
\frac1{V_n^2|\delta_n|}
=\frac{\beta_{n+1}G_n^2N^{u_{n+1}-2u_n}}
{|\alpha_{n+1}|D_n^2}.
\]
Writing $B_n=\den(\zeta_n)$, reduction can only decrease the denominator,
so
\begin{equation}\label{eq:coordinatecost}
B_n\le|\alpha_{n+1}|D_n^2N^{(2u_n-u_{n+1})_+}.
\end{equation}
The coefficient cost is $o(u_n)$ by \eqref{eq:finiteR}, and
$(2u_n-u_{n+1})_+=o(u_n)$ by the lower ratio hypothesis. Thus
\begin{equation}\label{eq:Bsmall}
\log B_n=o(M_n).
\end{equation}
This includes the steps with $u_{n+1}<2u_n$.
Finally,
\begin{equation}\label{eq:Lscales}
L_n=t_{n+1}+o(t_n),\qquad
\liminf\frac{L_n}{M_n}\ge2,\qquad
\limsup\frac{L_n}{M_n}=F.
\end{equation}
So the negative part of $L_n-2M_n$ is $o(M_n)$.

\section{A uniform finite profile bound}\label{sec:uniform}

In this section $F<\infty$; the case $F=+\infty$ follows from
\eqref{eq:lowerbound}.
Fix
\begin{equation}\label{eq:kappa}
\kappa>F-2.
\end{equation}
In particular, $\kappa>0$, also when $F=2$.
Put $C=\log6$. There are $n_0$ and $A\ge\log6$, depending on
$\kappa$ and the series, such that every logarithmic gap $a\to b$ of every $x_n$,
$n\ge n_0$, satisfies
\begin{equation}\label{eq:uniformprofile}
b-a\le\kappa a+A+C(n-n_0).
\end{equation}

Put $E_n=\log(4B_n)$. By \eqref{eq:Bsmall}--\eqref{eq:Lscales}
there is a sequence $\eta_n\to0$ such that
\begin{equation}\label{eq:tolerance}
\begin{gathered}
-\eta_nM_n\le L_n-2M_n\le(F-2+\eta_n)M_n,\\
E_n\le\eta_nM_n,\qquad \log(3/2)\le\eta_nM_n.
\end{gathered}
\end{equation}
Choose $n_0$ so that $V_n\ge2$ and
\begin{equation}\label{eq:tolerancechoice}
\eta_n\le\min\left\{\frac{\kappa-(F-2)}2,
\frac{\kappa}{16(1+\kappa)},
\frac{\kappa^2}{32(1+\kappa)}\right\}
\end{equation}
for $n\ge n_0$. Increase $A$, if necessary, to cover the finitely
many gaps of $x_{n_0}$.

Induct from $n$ to $n+1$ with Lemma~\ref{lem:perturbation}.
A bounded gap has width at most $\log6$. An inherited gap has the
same left endpoint as an old gap and width at most its width plus
$\log6=C$. Both satisfy \eqref{eq:uniformprofile} at index $n+1$.
A central gap begins at $a=M_n$ and has width at most
\[
L_n-2M_n\le(F-2+\eta_n)M_n\le\kappa a.
\]
For the short alternative,
\[
g\le2\eta_nM_n,\qquad a\ge(1-2\eta_n)M_n,
\]
and \eqref{eq:tolerancechoice} gives $g\le\kappa a$.

For a reflected gap, let $c\to d$ be its old gap and $g_0=d-c$.
Equations~\eqref{eq:widthandgain} and \eqref{eq:tolerance} give
\begin{equation}\label{eq:reflectioninduction}
g\le g_0+2\eta_nM_n,\qquad
a-c\ge g_0-2\eta_nM_n,\qquad
a\ge(1-2\eta_n)M_n.
\end{equation}
If $g_0<\kappa M_n/4$, then
\[
g\le(\kappa/4+2\eta_n)M_n
\le\kappa(1-2\eta_n)M_n\le\kappa a.
\]
If $g_0\ge\kappa M_n/4$, then
\begin{align*}
g-\kappa a
&\le(g_0-\kappa c)-\kappa g_0+(2+2\kappa)\eta_nM_n\\
&\le(g_0-\kappa c)-\frac{\kappa^2}{4}M_n
+\frac{\kappa^2}{16}M_n
\le g_0-\kappa c.
\end{align*}
Thus the bound for $c\to d$ gives the bound for the new gap without
loss, and \eqref{eq:uniformprofile} follows.

\section{Transfer to the infinite continued fraction}\label{sec:transfer}

Let $p/q$ be a convergent of $x$ with next denominator $q^+$, and put
\[
a=\log q,\qquad b=\log q^+,\qquad \varepsilon=|x-p/q|.
\]
Let $q\to\infty$. Gaps with $q^+/q\le6$ are bounded, and their widths
divided by the left endpoint tend to zero. Thus suppose $q^+/q>6$.
Choose the least $n\ge n_0$ for which
\begin{equation}\label{eq:chosentruncation}
t_n\ge4(a+b+1).
\end{equation}
Eventual geometric lower growth gives
\begin{equation}\label{eq:indexbound}
n=O(\log(a+b+2)).
\end{equation}
Indeed, $t_n\ge c_0\rho^n$ eventually for some $c_0>0$, and minimality gives
$t_{n-1}<4(a+b+1)$ once $n>n_0$.

The selected $n$ tends to infinity with $q$, since $a+b\to\infty$
and $t_n\ge4(a+b+1)$.
By \eqref{eq:reduceddenominator}--\eqref{eq:tail}, for all
sufficiently large selected indices,
\begin{equation}\label{eq:transfererror}
V_n\ge e^{t_n/2}>q,\qquad |x-x_n|\le e^{-t_n}.
\end{equation}
On the other hand,
\[
\varepsilon\ge\frac1{q(q^++q)}\ge\tfrac12e^{-(a+b)}.
\]
Thus $|x-x_n|<\varepsilon/2$. The inequality $V_n>q$ ensures
$p/q\ne x_n$, and
\begin{equation}\label{eq:transferlegendre}
0<\left|x_n-\frac pq\right|
\le\frac32\varepsilon\le\frac3{2qq^+}<\frac1{2q^2}.
\end{equation}
Legendre's criterion makes $p/q$ a nonterminal convergent of $x_n$.
Let $S$ be the next denominator after $q$ in the continued fraction of
$x_n$. Then
\eqref{eq:CFerror} and \eqref{eq:transferlegendre} give
\[
S\ge\tfrac23q^+-q\ge q^+/6.
\]
By \eqref{eq:uniformprofile} applied to this finite gap,
\[
b-a\le\log S-\log q+\log6
\le\kappa a+A+C(n-n_0)+\log6.
\]
By \eqref{eq:indexbound},
\begin{equation}\label{eq:infiniteprofile}
b-a\le\kappa a+O(\log(a+b+2)).
\end{equation}
The implied constants depend only on $\kappa$ and the series.

If $b\ge2(1+\kappa)a$, then \eqref{eq:infiniteprofile} gives
$b/2\le O(\log(2b+2))$, which fails for large $b$. Hence $b=O(a)$.
The term $C(n-n_0)$ in \eqref{eq:uniformprofile} is then $O(\log a)$,
and the error in \eqref{eq:infiniteprofile} is $O(\log a)=o(a)$.
Consequently,
\begin{equation}\label{eq:limitinggap}
\limsup\frac{b-a}{a}\le\kappa.
\end{equation}

\section{Completion of the proof}

For an irrational $x$ with convergent denominators $q_m$,
\eqref{eq:CFerror} gives
\[
-\frac{\log|x-p_m/q_m|}{\log q_m}
=1+\frac{\log q_{m+1}}{\log q_m}+o(1).
\]
By Legendre's criterion every reduced approximation with
$|x-p/q|<1/(2q^2)$ is a convergent; so is every reduced approximation of
fixed exponent greater than $2$ once $q$ is large.
In the definition of $\mu$, a pair $(p,q)$ need not be reduced. Reducing it does not increase the denominator, so the same inequality holds. A
fixed reduced fraction occurs only finitely often among approximations of a
fixed positive exponent, because its error is nonzero. In a bounded interval
there are only finitely many reduced fractions with bounded denominator.
Thus infinitely many such approximations have unbounded reduced
denominators, and the preceding Legendre argument applies when the
exponent is greater than $2$.
Therefore
\begin{equation}\label{eq:exponentidentity}
\mu(x)=1+\limsup_m\frac{\log q_{m+1}}{\log q_m}
=2+\limsup_m\frac{\log(q_{m+1}/q_m)}{\log q_m}.
\end{equation}
This holds also when the limsup is infinite.
Equation~\eqref{eq:limitinggap} gives $\mu(x)\le2+\kappa$.
Letting $\kappa$ decrease to $F-2$ and using
\eqref{eq:lowerbound} proves \eqref{eq:main} for finite $F$.
The infinite case is in Section~\ref{sec:arithmetic}, which completes
the proof of Theorem~\ref{thm:main}.\qed

\section{Height sharpness and scope}

The coefficient condition cannot be replaced by
$\log H(\lambda_j)\le C_*u_j$ with any prescribed constant $C_*>0$.
The first construction below shows this. The second shows that, at one
fixed linear height scale, the sum can be rational or have
infinite irrationality exponent.
In this section index the exponents by $j\ge0$ and take $u_j=2^j$.
Write $\ell=\log N$.

Fix $0<\theta<1/3$, and set
\[
m_j=\lfloor\theta2^j\rfloor,\qquad
\lambda_j=N^{(-1)^jm_j},\qquad e_j=2^j-(-1)^jm_j.
\]
These positive rational weights give
\[
\log H(\lambda_j)=m_j\ell=\theta\ell2^j+O(1),\qquad
x=\sum_{j\ge0}N^{-e_j}.
\]
The exponents $e_j$ are strictly increasing. For odd $j$ their
difference is at least $(1-3\theta)2^j>0$; for even $j$ it is
$2^j+m_j+m_{j+1}>0$. Their gaps tend to infinity, so the base-$N$
expansion is not eventually periodic and $x$ is irrational.
The truncation $p_j/q_j$ through index $j$ has reduced denominator
$q_j=N^{e_j}$, since its numerator is $1$ modulo $N$, and
\[
0<x-p_j/q_j\le\frac{N^{-e_{j+1}}}{1-N^{-1}}
\le2N^{-e_{j+1}}.
\]
Along even indices $e_{j+1}/e_j\to2(1+\theta)/(1-\theta)$, whence
\begin{equation}\label{eq:sharpheight}
\mu(x)\ge\frac{2(1+\theta)}{1-\theta}>2.
\end{equation}
Choosing $0<\theta<\min\{1/3,C_*/\ell\}$ makes
$\log H(\lambda_j)\le C_*2^j$ for every $j$.

At the fixed linear height scale below, positive integer weights can give a
rational sum. Set
\[
\lambda_0=N-1,\qquad \lambda_j=N^{2^{j-1}}-1\quad(j\ge1).
\]
Then $\lambda_jN^{-2^j}=N^{-2^{j-1}}-N^{-2^j}$ for $j\ge1$, so
\[
\sum_{j=0}^m\lambda_jN^{-2^j}=1-N^{-2^m}\longrightarrow1,
\qquad \log H(\lambda_j)\sim\tfrac12\ell2^j.
\]
At the same scale, replace the later weights by
\[
\widetilde\lambda_j=N^{2^{j-1}}-1+
\boldsymbol1_{\{j=r^2\text{ for some }r\ge1\}}\quad(j\ge1),
\qquad\widetilde\lambda_0=N-1.
\]
They remain positive integers, with logarithmic height
asymptotic to $\tfrac12\ell2^j$. Their sum is
\[
\widetilde x=1+\xi_N,\qquad
\xi_N=\sum_{r\ge1}N^{-2^{r^2}}.
\]
The digit expansion shows that $\xi_N$ is irrational. The truncation
through $r=J$ has reduced denominator
$q_J=N^{2^{J^2}}$, since its numerator is one modulo $N$, and error
\[
0<\xi_N-\sum_{r=1}^JN^{-2^{r^2}}
\le2N^{-2^{(J+1)^2}}=2q_J^{-2^{2J+1}}.
\]
Adding one to each approximant preserves its reduced denominator and its
approximation error. Thus $\mu(\widetilde x)=+\infty$.

The height condition is sufficient but not necessary for a given series. For example,
if $b\ge2$ is an integer, the weights
$\lambda_j=b^{-2^j}$ have logarithmic height $(\log b)2^j$,
but their sum is $\sum_{j\ge0}(Nb)^{-2^j}$, which has exponent $2$
by Theorem~\ref{thm:main} with base $Nb$ and unit weights.

Theorem~\ref{thm:main} does not cover
$\liminf u_{j+1}/u_j<2$. The examples above keep the exponent ratios
equal to $2$ and vary only the coefficient heights.
Problem~10.40 in \cite[p.~220]{Bu12} asks for the irrationality exponent
of binary lacunary sums $\sum_j2^{-n_j}$ for an arbitrary increasing
lacunary exponent sequence $(n_j)$. For $N=2$, Theorem~\ref{thm:main} gives the exponent when
$\liminf n_{j+1}/n_j\ge2$, including ratios tending to $2$ from below,
and allows small rational weights. The case with eventual ratio at least
$2$ is recorded in \cite{BKS} for unit coefficients, with attribution
there to Bugeaud's earlier work, and for coefficients $\pm1$ it also
follows from \cite[Theorem~1]{DKS}. Under the same height hypothesis, the
subcase $2<\liminf n_{j+1}/n_j\le\limsup n_{j+1}/n_j<\infty$ also follows
from \cite[Theorem~1]{DKS}.

\section*{Use of AI tools}

The author devised the main argument. Codex (OpenAI, GPT-6 Astra) carried
out exact checks of the continued-fraction identities and searched the literature.
Claude Code (Anthropic, Claude Opus 5.5) reviewed the arguments and
references. The author takes responsibility for the proof and references.

\begin{thebibliography}{99}

\bibitem{AB}
M.~Amou and Y.~Bugeaud, \emph{Exponents of Diophantine approximation
and expansions in integer bases}, J. Lond. Math. Soc. (2)
\textbf{81} (2010), no.~2, 297--316,
\url{https://doi.org/10.1112/jlms/jdp070}.

\bibitem{Bu08}
Y.~Bugeaud, \emph{Diophantine approximation and Cantor sets},
Math. Ann. \textbf{341} (2008), no.~3, 677--684,
\url{https://doi.org/10.1007/s00208-008-0209-4}.

\bibitem{Bu12}
Y.~Bugeaud, \emph{Distribution Modulo One and Diophantine Approximation},
Cambridge Tracts in Mathematics, vol.~193, Cambridge University Press,
Cambridge, 2012, \url{https://doi.org/10.1017/CBO9781139017732}.

\bibitem{BKS}
Y.~Bugeaud, D.~Krieger, and J.~Shallit,
\emph{Morphic and automatic words: maximal blocks and Diophantine
approximation}, Acta Arith. \textbf{149} (2011), no.~2, 181--199,
\url{https://doi.org/10.4064/aa149-2-7}.

\bibitem{DKS}
D.~Duverney, T.~Kurosawa, and I.~Shiokawa,
\emph{Irrationality exponents of certain fast converging series of
rational numbers}, Tsukuba J. Math. \textbf{44} (2020), no.~2,
235--250, \url{https://doi.org/10.21099/tkbjm/20204402235}.

\bibitem{MF}
M.~Mend\`es France, \emph{Sur les fractions continues limit\'ees},
Acta Arith. \textbf{23} (1973), no.~2, 207--215,
\url{https://doi.org/10.4064/aa-23-2-207-215}.

\bibitem{vPS}
A.~J.~van der Poorten and J.~Shallit, \emph{Folded continued fractions},
J. Number Theory \textbf{40} (1992), no.~2, 237--250,
\url{https://doi.org/10.1016/0022-314X(92)90042-N}.

\bibitem{Sh79}
J.~O.~Shallit, \emph{Simple continued fractions for some irrational
numbers}, J. Number Theory \textbf{11} (1979), no.~2, 209--217,
\url{https://doi.org/10.1016/0022-314X(79)90040-4}.

\bibitem{Sh82}
J.~O.~Shallit, \emph{Simple continued fractions for some irrational
numbers, II}, J. Number Theory \textbf{14} (1982), no.~2, 228--231,
\url{https://doi.org/10.1016/0022-314X(82)90047-6}.

\bibitem{Zh}
R.~Zhang, \emph{The irrationality exponent of small-height binary
lacunary series is 2}, Zenodo concept record (2026),
\url{https://doi.org/10.5281/zenodo.24040635}.

\end{thebibliography}

\end{document}
